\chapter{导数在高中数学最优化问题中的应用}

为使学生充分认识到导数在解决最优化问题中的功能与价值，教师要更关注与其相关的最大利润、最节约材料、用时最少等问题。


\section{最大利润问题}

\begin{example}
  一家工厂生产不同等级的衬衣，按品质分成12个等级，第一等级的衬衣（也就是最低级的衬衣），它的利润是十块钱，如果提高一件衬衣的等级，利润会增加五块钱，而在同样的时间里生产的量会少五件。假设在相同时间内，第一等级的衬衫能生产100件。问：同一时间内，生产何种等级的衬衣，其整体盈利最高？最大利润有多少？\cite{张双林2013导数在实际生活中的最优化应用}
\end{example}

\begin{analysis}
  最大利润问题往往归结为高中数学中求函数的最值问题，利用求导法可以解决这类问题，但是需要注意极值是否包含在函数$f(x)$的定义域内。
\end{analysis}

\begin{proof}[解析]
  假设在同一时间内, 第 ${x}\left(x \in \mathbb{N}^{*}, 1 \leq x \leq 12\right)$ 等级的祇衫生产后获利润 ${y}$ 是 最大的。根据题意, 得
  \[
  y=[10+5(x-1)][100-5(x-1)]=25(x+1)(21-x).\] 

  对 $y$ 求导数, 得 $y^{\prime}=25(21-x)-25(x+1)=50(10-x)$,
  令 $y^{\prime}=0$, 解得 $x=10$. 因 $x=10 \in[1,12]$, 所以 $y$ 只有一个极值点且 $x=10$ 时, $y=3025$.
  又因为 $x=1$ 时, $y=1000 ; x=12$ 时, $y=2925$. 所以3025是原函数的最大值。 因此 在等量时间内, 利润最高可达3025元.此时生产的为第十等级衬衫.
\end{proof}

\begin{review}
  对于高中数学最优化问题，假如目标函数是二次及以上的多项式函数、幂函数、对数函数和指数函数以及它们的复合函数等等，则可通过求导法求出最优解。因此，导数求解问题扩大了导数知识在高中数学最优化问题中的应用范围。
\end{review}



\section{用料最省问题}

\begin{example}
  一家饮料工厂要制造一种圆柱形的金属罐头，在确定了体积一定后，应该如何选择其高、底和半径，以减少使用的材料？
\end{example}


\begin{proof}[解析]
  设圆柱的高为 $h$, 底半径为 $R$, 则易拉罐表面积 $S(R)=2 \uppi R h+2 \uppi R^{2}$ ，由 $V=$ $\uppi R^{2} h$, 得 $h=\frac{V}{\uppi R^{2}}$, 则 $S(R)=2 \uppi R \frac{V}{\uppi R^{2}}+2 \uppi R^{2}=\frac{2 V}{R}+2 \uppi R^{2}$, 令 $S^{\prime}(R)=-\frac{2 V}{R^{2}}+4 \uppi R=0$  , 
  解得 $R=\sqrt[3]{\frac{V}{2 \uppi}}$ ,从而 $h=\frac{V}{\uppi R^{2}}=\frac{V}{\uppi(\sqrt[3]{\frac{V}{2 \uppi}})^{2}}=\sqrt[3]{\frac{4 V}{\uppi}}=2 \sqrt[3]{\frac{V}{\uppi}}$ 即 $h=2 R$.
  当 $h<2 R$ 时, $S^{\prime}(R)<0$; 当 $h>2 R$ 时, $S^{\prime}(R)>0$; 因此它是原函数$S(R)$的极小值，而且是唯一的极值，所以它是最小值，此时所用材料最少。
\end{proof}


\section{追击用时最少问题}

\begin{example}
  张三下了车, 发现自己的东西落在了大巴车上, 这时大巴车已经停在了一个十字 路口, 张三用 $6 \unit{m\per s}$ 的速度去追前面的大巴车, 在距离大巴车 $25 \unit{m}$ 的时候, 红灯变成 了绿灯, 大巴车以 ${a}=1 \unit{m \per s^2}$ 的加速度开始行驶, 问张三能不能赶上大巴车? 如果赶不上, 请求出张三和大巴车之间的最小距离\cite{张双林2013导数在实际生活中的最优化应用}。
\end{example}

\begin{analysis}
  运用位移公式和它们间的距离关系, 构造二次函数模型进行求导．
\end{analysis}

\begin{proof}[解析]
  假如 ${t}$ 秒后张三能追上大巴车, 此时人的位移为 $S_{1}=6 t$ 且大巴车的位移为 $S_{2}=\frac{1}{2} {at}^{2}=\frac{t^{2}}{2}$, 根据题意得, $S_{1}-S_{2}=6 t-\frac{t^{2}}{2}=25$, 化简得 $t^{2}-12 t+50=0$, 由于判别 式 $\Delta=12^{2}-200<0$, 所以此方程无解, 因此张三无法追上大巴车。

  那么我们接下来就是求经过时间 ${t}$ 秒后张三与大巴车之间的距离 $f(t)=S_{2}+25-S_{1}=\frac{t^{2}}{2}+25-6 t$, 求导得 $f^{\prime}(t)=t-6$, 令 $f^{\prime}(t)=0$, 得 $t=6$. 当 ${t}<6$ 时, ${f}^{\prime}({t})<0$; 当 $t>6$ 时, ${f}^{\prime}(t)>0$.

  因此, 当 ${t}=6$ 时, 张三与大巴车的距离最近, 最小距离为 7 米。
\end{proof}

\begin{review}
  该问题实际上是一个求二次函数最小值的问题，除了求导法外，还可以用配方法、顶点坐标法、图像法等求二次函数最值。
\end{review}